\documentclass[11pt]{article}

\usepackage[margin=1in]{geometry}
\usepackage[T1]{fontenc}
\usepackage{lmodern}
\usepackage{amsmath}
\usepackage{amssymb}
\usepackage{booktabs}
\usepackage[colorlinks=true,linkcolor=blue]{hyperref}
\usepackage{cleveref}

\title{Advanced Display Mathematics: A Stress Test for Structured Equation Conversion}
\author{Test Corpus Author}
\date{\today}

\begin{document}

\maketitle

\begin{abstract}
This note exercises a range of advanced \texttt{amsmath} and \texttt{amssymb}
constructions that frequently appear in analysis and combinatorics. The
objective is to present these features in genuine mathematical prose, with
numbered equations and cross-references, so that any document-conversion
pipeline targeting Office Math Markup Language must reproduce subscript stacks,
decorated arrows, braces, multi-line displays, boxed results, and modular
arithmetic notation.
\end{abstract}

\section{Sums over restricted index sets}
\label{sec:sums}

A common pattern in enumerative combinatorics restricts a summation to index
pairs satisfying several simultaneous conditions. We write such constraints
beneath the summation sign using a stacked subscript. For nonnegative integers
$i$ and $j$ bounded by $n$, consider the bilinear quantity
\begin{equation}
\label{eq:substack}
S_n \;=\; \sum_{\substack{0 \le i \le n \\ 0 \le j \le n \\ i + j = n}}
   \binom{n}{i}\, \binom{n}{j},
\end{equation}
where the inner factors are ordinary binomial coefficients. Because the
constraint $i + j = n$ collapses the double sum to a single free index,
\cref{eq:substack} evaluates in closed form to $\binom{2n}{n}$ by the
Vandermonde convolution. The notation $\binom{n}{i}$ should render as a genuine
two-row binomial coefficient rather than a fraction.

\section{Decorated relations and limits}
\label{sec:decorated}

When two operators act in sequence it is often clarifying to annotate a relation
symbol with the transformation responsible for it. Using \verb|\overset| and
\verb|\underset| we display the data of a map together with its name:
\begin{equation}
\label{eq:overset}
f(x) \;\overset{\text{def}}{=}\; \lim_{t \to 0^{+}} \frac{g(x+t) - g(x)}{t},
\qquad
a_n \;\underset{n \to \infty}{\longrightarrow}\; L .
\end{equation}
The annotation above the equality sign in \cref{eq:overset} marks it as a
definition, while the label beneath the arrow records the limiting regime.

A directed construction such as a chain complex or a sequence of refinements is
naturally expressed with labelled extensible arrows. The morphisms below carry
their names above and their degrees below:
\begin{equation}
\label{eq:arrows}
A \xrightarrow{\;\varphi\;} B \xleftarrow[\psi]{\quad} C,
\qquad
X_0 \xrightarrow[\text{stage } 1]{\;\pi_1\;} X_1 .
\end{equation}
Here \cref{eq:arrows} pairs a right-pointing arrow decorated with $\varphi$ and a
left-pointing arrow decorated beneath by $\psi$, each of which must retain its
direction and its annotation after conversion.

\section{Braced groupings}
\label{sec:braces}

Over- and under-braces are indispensable when a formula must be read as a few
meaningful blocks. Grouping the terms of a finite geometric expansion makes the
recursive structure visible:
\begin{equation}
\label{eq:braces}
\overbrace{1 + x + x^{2} + \dots + x^{k}}^{\text{partial sum } P_k}
\;=\;
\underbrace{\frac{1 - x^{k+1}}{1 - x}}_{\text{valid for } x \neq 1}.
\end{equation}
The superscripted label on the overbrace and the subscripted side condition on
the underbrace in \cref{eq:braces} are essential to the reading and must survive
as attached annotations rather than detached text.

\section{A long display split across lines}
\label{sec:multline}

Some identities are too wide for a single typeset line and are broken with the
\texttt{multline} environment, which sets the first part flush left and the last
part flush right. Expanding a triple product gives
\begin{multline}
\label{eq:multline}
(a + b + c)^{3}
= a^{3} + b^{3} + c^{3} \\
+ 3a^{2}b + 3a^{2}c + 3b^{2}a + 3b^{2}c + 3c^{2}a + 3c^{2}b \\
+ 6abc ,
\end{multline}
and the grouping of \cref{eq:multline} into a leading cube-sum line, a middle
line of mixed quadratic terms, and a trailing symmetric term should be preserved
in the converted output.

\section{Boxed conclusions, congruences, and aligned derivations}
\label{sec:align}

To emphasise a final result we enclose it in a box. Fermat's little theorem, for
a prime $p$ and an integer $a$ not divisible by $p$, asserts
\begin{equation}
\label{eq:boxed}
\boxed{\,a^{\,p-1} \equiv 1 \pmod{p}\,},
\end{equation}
where the modular qualifier produced by \verb|\pmod| sits in parentheses to the
right of the congruence. The boxed display in \cref{eq:boxed} must remain a
single framed unit.

The following aligned derivation carries a custom equation marker supplied by
\verb|\tag*|, which suppresses the automatic number while still labelling the
line for reference:
\begin{align}
\label{eq:star}
\sum_{k=0}^{n} k
&= 0 + 1 + 2 + \dots + n \notag \\
&= \frac{n(n+1)}{2}. \tag*{($\ast$)}
\end{align}
We may refer back to the starred line via its tag using \eqref{eq:star}, which
should display the literal marker $(\ast)$ in parentheses rather than a serial
number.

\section{Piecewise definitions and named operators}
\label{sec:cases}

Functions defined by cases use the \texttt{cases} environment, which aligns the
branch values against their guards. The sign function is the canonical example,
\begin{equation}
\label{eq:cases}
\operatorname{sgn}(x) =
\begin{cases}
  +1 & \text{if } x > 0, \\
  \phantom{+}0 & \text{if } x = 0, \\
  -1 & \text{if } x < 0,
\end{cases}
\end{equation}
where \verb|\operatorname| typesets \(\operatorname{sgn}\) in upright roman as a
named operator. A further example combining a named operator with modular
reduction is the residue map
\begin{equation}
\label{eq:opmod}
\operatorname{res}_m(x) \;=\; x \bmod m
\;=\; x - m \left\lfloor \frac{x}{m} \right\rfloor ,
\end{equation}
so that $\operatorname{res}_m(x) \equiv x \pmod{m}$. The piecewise alignment in
\cref{eq:cases} and the operator names in \cref{eq:cases,eq:opmod} are the
features under test.

\section{Summary}
\label{sec:summary}

The displays of \cref{sec:sums,sec:decorated,sec:braces,sec:multline,sec:align,sec:cases}
collectively exercise restricted-index summation, decorated relations and arrows,
braced groupings, a multi-line split, a boxed congruence with a modular
qualifier, a starred custom tag recalled by \eqref{eq:star}, a piecewise
definition, and named operators. A faithful conversion preserves the structure
of each construction
together with its attached annotations.

\end{document}
